% Standalone section: theorem/proof environments and amsmath required.
% Bibliography keys: MagnusKoszul, LalondeRam (entries suggested at end).

\section{The intersection over all lift choices}
\label{sec:all-lift-intersection}

We now determine the exact strength of requiring the field inclusion
for every unrestricted choice of lifts.  No invariance under
$\Delta$ or $\Gamma$ is imposed on the individual normal closures
in this section.

\begin{theorem}\label{thm:all-lift-intersection}
Let $P$ be a free pro-$p$ group on a finite or countably infinite
convergent free basis, of rank at least two, and put
\[
 F=P*F(t_1,\ldots,t_r),\qquad r<\infty.
\]
For $\boldsymbol w=(w_1,\ldots,w_r)\in[P,P]^r$, let
$q_{\boldsymbol w}:F\to P$ fix $P$ and send $t_j$ to $w_j$.
Then
\[
 \bigcap_{\boldsymbol w\in[P,P]^r}\ker q_{\boldsymbol w}=1,
 \qquad
 \ker q_{\boldsymbol w}
   =\overline{\langle t_jw_j^{-1}:1\leq j\leq r\rangle}^{\rm normal}_F.
\]
Each $t_jw_j^{-1}$ has the same image as $t_j$ in $F^{\rm ab}$.
For every nonidentity element the proof supplies a separating
homomorphism to a finite $p$-group which factors through one of
these retractions.
\end{theorem}

\begin{proof}
The assertion for $r=0$ is immediate.  First assume that $P$ has
finite rank $d\geq2$, with generators $x_1,\ldots,x_d$.
Use the Magnus embedding
\[
 \mu:F\hookrightarrow
 1+(X_1,\ldots,X_d,T_1,\ldots,T_r)
 \subset\mathbb F_p\langle\!\langle X_1,\ldots,X_d,T_1,\ldots,T_r
                       \rangle\!\rangle,
\]
where $x_i\mapsto1+X_i$ and $t_j\mapsto1+T_j$.
The completed group-algebra identification and its injectivity
are recalled in \cite[\S3.1, equation~(3.3), p.~8]{MagnusKoszul}.

Fix $h\ne1$.  If the retraction $t_j\mapsto1$ does not kill $h$,
it already separates $h$.  Otherwise every monomial of
$\mu(h)-1$ contains a $T$-letter.  Let $k\geq1$ be the smallest
number of $T$-letters occurring in a monomial with nonzero
coefficient.  Among monomials with exactly $k$ such letters,
let $\ell$ be the smallest number of $X$-letters.
Their sum is a nonzero polynomial $Q$ over $\mathbb F_p$.
Every word of $Q$ has exactly $k$ $T$-letters and $\ell$
$X$-letters.

Choose two of the retained variables and call them $X,Y$,
ordered by $X<Y$.  Choose an integer $D>\ell$ and, for
$1\leq j\leq r$, put
\[
 u_j=X Y^D X Y^{D+j} X Y^{D+r+1-j}.
\]
All these words have the same length
$L=3D+r+4>\ell$.
Each is a Lyndon word.  Indeed a rotation beginning with $Y$
is larger.  A rotation beginning at either later occurrence
of $X$ has an initial $Y$-run longer than $D$, so it is larger
than the rotation beginning at the first $X$.
Let $W_j$ be its integral standard Lyndon bracketing reduced
modulo $p$.  The triangular expansion is
\[
 W_j=u_j+\sum_{v>u_j,\ |v|=L}a_vv.
\]
In particular the coefficient of $u_j$ is one in every
characteristic; see \cite[\S1, equation~(1.2)]{LalondeRam}.
Replacing the Lie brackets by group commutators gives
$w_j\in[P,P]$ whose Magnus expansion begins with $W_j$ in
degree $L$.

We check the word substitution $T_j\mapsto u_j$ precisely.
In any substituted word from $Q$, the occurrences of the
signature
\[
                             X Y^D X
\]
are exactly the beginnings of the inserted markers $u_j$.
Inside a marker its first $Y$-run has length $D$ and its
other two $Y$-runs have lengths strictly larger than $D$.
Between markers the total number of original $X_i$-letters
is $\ell<D$, so those letters cannot form a $Y$-run of
length $D$.  A final $Y$-run of a marker can merge with
following original $Y$-letters, but its length is already
larger than $D$ and remains so.  The initial $X$ of the
next marker prevents merger with its first $Y$-run.
These observations also rule out a signature crossing a
marker boundary.

The signature therefore determines every marker's starting
position.  The length of its middle $Y$-run is $D+j$, so it
determines $j$; the common length $L$ then determines its
ending position, even when its final run merges with
following original letters.  Removing the markers recovers
the original word of $Q$.  Thus distinct words of $Q$ have
distinct substituted words.

Choose the lexicographically least substituted word among
those with a nonzero coefficient in $Q$.
The triangular expansions of the $W_j$ show that it occurs
with that same coefficient in $Q(X_1,\ldots,X_d,W_1,\ldots,W_r)$.
Every other contribution is larger: within one substitution
this follows from the equal lengths of the terms of $W_j$,
and between different words it follows from the choice of
the least substituted word.  Therefore this polynomial is
nonzero.

In the full substituted series, its degree is $kL+\ell$.
A monomial with more than $k$ $T$-letters has degree at least
$(k+1)L>kL+\ell$.  A monomial with exactly $k$ $T$-letters
and more than $\ell$ $X$-letters has larger degree as well.
The higher-degree terms of $\mu(w_j)-1$ also contribute
only in larger degree.  Consequently
\[
 \mu(q_{\boldsymbol w}(h))-1
\]
has the nonzero initial polynomial just exhibited, so
$q_{\boldsymbol w}(h)\ne1$.
Truncating the noncommutative series algebra after degree
$kL+\ell$ gives a finite-dimensional nilpotent augmentation
ideal $J$ over $\mathbb F_p$.
The finite group $1+J$ is a $p$-group and still detects this
nonzero coefficient.  This is the promised finite quotient.

For countably infinite rank, a homomorphism from $F$ to a
finite group kills all but finitely many members of its
convergent free basis.  Choose such a homomorphism detecting
$h$, and retain the finitely many $x_i$ it does not kill,
adding two retained generators when necessary.  The
retraction which kills the other $x_i$ and keeps all $t_j$
sends $h$ to a nonidentity element of a finite-rank free
pro-$p$ group.  Apply the preceding construction there.
Its $w_j$ use only the retained generators, so the finite-rank
retraction commutes with the substitutions.  The resulting
$q_{\boldsymbol w}$ on the original $F$ still detects $h$.

Finally the displayed normal-closure formula follows from
the universal property of the free pro-$p$ product: adjoining
the relations $t_j=w_j$ gives exactly $P$, with the indicated
quotient map.  The abelian image assertion follows from
$w_j\in[P,P]$.
\end{proof}

\begin{corollary}\label{cor:all-lifts-faithfulness}
Let $A$ be a finite-rank direct summand of the abelianization
of a countably based free pro-$p$ group $F$, and suppose the
complementary free rank is at least two.  Fix a
$\mathbb Z_p$-basis of $A$.  The intersection of the closed
normal closures of all tuples lifting this basis is trivial.
For any homomorphism $\rho:F\to G$, therefore,
\[
 \ker\rho\subseteq N\ \text{for every such lift normal closure }N
       \quad\Longleftrightarrow\quad \ker\rho=1.
\]
\end{corollary}

\begin{proof}
A finite tuple lifting the basis of $A$ has independent
Frattini images and extends to a convergent free basis, by
the finite-coordinate and finite-tail argument in
Theorem~\ref{thm:component-no-galois-p37}.
Write the resulting free-product decomposition as in
Theorem~\ref{thm:all-lift-intersection}.  The tuples
$t_jw_j^{-1}$ in that theorem are among the tuples under
consideration.  Their normal closures already have trivial
intersection.  The equivalence is now the literal
intersection of the displayed subgroup containments.
\end{proof}

In the cyclotomic application, the omitted sum of even
components is finite free over $\mathbb Z_p$ and is a direct
summand of $H^{\rm ab}$.  The complementary free rank is
infinite because the odd part has positive $\Lambda$-rank.
Consequently a field inclusion required for every unrestricted
lift choice,
\[
                   \Omega_p^N\subseteq\Omega_p^*,
\]
is, through Galois correspondence and
Corollary~\ref{cor:all-lifts-faithfulness}, equivalent to the
full equality $\Omega_p=\Omega_p^*$.
This statement concerns the quantifier over all lift choices;
it adds no Galois-invariance assertion for an individual $N$.

% Suggested bibliography entries:
% \bibitem{MagnusKoszul}
% J. Minac, F. W. Pasini, C. Quadrelli, N. D. Tan,
% Mild pro-p groups and the Koszulity conjectures,
% Expositiones Math. 40 (2022), 432--455; arXiv:2106.03675v2,
% section 3.1, equation (3.3), printed p.8.
% \bibitem{LalondeRam}
% P. Lalonde and A. Ram, Standard Lyndon bases of Lie algebras
% and enveloping algebras, Trans. Amer. Math. Soc. 347 (1995),
% 1821--1830, section 1, equation (1.2).
